\documentclass[11pt]{article}
\usepackage{amsmath,amssymb,amsthm,geometry}\geometry{margin=1in}
\newtheorem{theorem}{Theorem}
\title{A naturality boundary for marginal-only dependence}
\date{15 November 2026}\begin{document}\maketitle
\paragraph{Status and level.} A self-contained finite-distribution boundary argument proposed for independent review. This is a semantic obstruction informing the desired general interaction structure, not that structure itself. No novelty claim is made. It says which information cannot be omitted, not that general shared dependence is impossible.

Let $D$ be the finite-support normalized-distribution functor on Set. Suppose a family of functions
\[
 c_{X,Y}:DX\times DY\longrightarrow D(X\times Y)
\]
preserves both marginals and is natural for \emph{independently chosen} deterministic maps $f:X\to X'$, $g:Y\to Y'$:
\[
 D(f\times g)c_{X,Y}(\mu,\nu)
   =c_{X',Y'}(Df\mu,Dg\nu).\tag{N}
\]
No continuity, measurability, associativity, symmetry or monad law for $c$ is assumed.

\begin{theorem}
The family $c$ is necessarily the ordinary independent product: $c_{X,Y}(\mu,\nu)(x,y)=\mu(x)\nu(y)$.
\end{theorem}
\begin{proof}
Write $b_p$ for the Bernoulli distribution on $\{0,1\}$ with mass $p$ at $1$. Let
\[
 F(p,q)=c_{2,2}(b_p,b_q)(1,1).
\]
It is nonnegative. For $p_1,p_2\ge0$ with $p_1+p_2\le1$, consider the distribution with masses $p_1,p_2,1-p_1-p_2$ on a three-point set. Pair it using $c$ with $b_q$. Naturality for the indicators of the first point, the second point, and their union gives
\[
 F(p_1+p_2,q)=F(p_1,q)+F(p_2,q).
\]
This is ordinary additivity of the \emph{one fixed three-by-two joint distribution}; it does not assume that $c$ is affine in either argument. Marginal preservation gives $F(1,q)=q$ and $F(0,q)=0$. Nonnegativity plus the displayed additivity makes $F(-,q)$ monotone. Integer subdivision therefore gives $F(m/n,q)=(m/n)q$, and rational upper/lower bounds prove $F(p,q)=pq$ for every real $p\in[0,1]$. No continuity premise was used.

For arbitrary finite $X,Y$ and any $x,y$, apply (N) to the singleton indicators $1_{\{x\}},1_{\{y\}}$. The $(1,1)$ mass becomes exactly the joint mass at $(x,y)$, so it equals $F(\mu(x),\nu(y))=\mu(x)\nu(y)$. For arbitrary sets and finite-support distributions, marginal preservation forces zero mass outside the product of the two finite supports; the same singleton-indicator argument applies. Empty carriers have no normalized input distributions and impose no exception. The product family itself has the required naturality and marginals.
\end{proof}

In particular no such family can also choose the diagonal coupling $D\Delta_X(\mu)$ whenever both inputs equal a non-Dirac $\mu$: the product has positive off-diagonal mass between two distinct positive atoms. This does not prohibit a diagonal operation on a single supplied distribution, or naturality restricted to the \emph{same} map in both coordinates. The natural single-distribution diagonal does not itself construct a coupling rule for arbitrary unequal marginals; the existence/classification of a binary family under only same-map naturality is left undecided here. The theorem requires the full independent pair of pure maps in (N), and its input data consist solely of the two marginals.

\section{Consequence for the categorical design problem}
A general dependence-preserving interaction must retain more than those marginals: for example an actual joint witness, a common boundary with supplied conditional data, or an admitted relation/interaction policy. Alternatively it can restrict the pure maps or weaken the requested naturality. These are genuine choices of structure, not consequences of arbitrary affine commutative monad axioms. The result explains why replacing a retained joint context by its component projections and then asking for a fully natural reconstruction would silently force independence in this semantic instance.

Conversely supplying an arbitrary pairwise coupling does not automatically supply a coherent triple interaction, conditional evolution or feedback. The outstanding general theorem must specify those extra data and derive its assembly laws from local hypotheses; it cannot simply assume the desired whole-system joint square. This note neither constructs that theorem nor identifies one particular stochastic realization with the general categorical core.
\end{document}
